\section{From the Dual Dirac Equation to each observable's prediction}
\label{sec:derivation-chain}

\noindent\textit{This section transcribes the derivation chain from
\texttt{Roadmapping/Equation\_Verification/Dual\_Relativistic\_Quantum\_Mechanics\_I.md}
(DRQM~I, \S\S II--III) and the per-PR Bethe--Salpeter section docs
(\texttt{01\_NonRelHydrogen.md}, \texttt{03\_FineStructure.md},
\texttt{05\_LambShift.md}, \texttt{06\_Hyperfine.md}). No new physics is
introduced; the chain is presented in the form the senior author actually
reads (PDF with LaTeX equations, not Mathematica notebook). Every
numbered DRQM~I equation from (II.1) through (III.22) that carries a
typeset form in the verification doc is surfaced below for fine-grain
review, in order. Equations DRQM~I records structurally rather than
typesetting (the intermediate (III.11)--(III.16) spherical-coordinate steps)
are noted honestly as such; their three load-bearing
results~(III.17-equivalent) are displayed.}

\subsection{The proper-time canonical form}

DRQM~I introduces three dual relativistic wave equations, all of the form
\begin{equation}
\label{eq:dual-K-form}
i\hbar\,\frac{\partial\Psi}{\partial\tau} \;=\; K\,\Psi,\qquad
K \;=\; \frac{H^2}{2 m c^2} \;+\; \frac{m c^2}{2},\qquad
H \in \{H_{\rm Dirac},\;H_{\rm sqrt(1)},\;H_{\rm sqrt(2)}\}.
\end{equation}
The three are distinguished by which $H$ is substituted into~\eqref{eq:dual-K-form}.

\subsection{The Dual Dirac Equation, Eq.~(II.1)}

With $H_{\rm Dirac} = c\,\boldsymbol\alpha\!\cdot\!\boldsymbol\pi
+ \beta m c^2 + V$ and $\boldsymbol\pi = \mathbf{p} - e\mathbf{A}/c$, the
Dual Dirac Equation is
\begin{equation}
\label{eq:II1-dual-dirac}
\begin{aligned}
i\hbar\,\frac{\partial\Psi}{\partial\tau} = \Bigg\{&\,
   \frac{\boldsymbol\pi^2}{2m}
 + \beta V
 + m c^2
 - \frac{e\hbar\,\boldsymbol\Sigma\!\cdot\!\mathbf{B}}{2 m c} \\
&{}+ \frac{V\,\boldsymbol\alpha\!\cdot\!\boldsymbol\pi}{m c}
 - \frac{i \hbar\,\boldsymbol\alpha\!\cdot\!\nabla V}{2 m c}
 + \frac{V^2}{2 m c^2}\,
\Bigg\}\Psi.
\end{aligned}
\end{equation}
\wolframcheck{\texttt{FullSimplify} of $K_{\rm Dirac} - (\text{II.1})$ returns
$0$ (DRQM~I~\S II.1 verified, 2026-05-11).}

\subsection{The cleanest form: Eq.~(II.3)}

Eq.~(II.3) (DRQM~I~\S II.3) takes
$H_{\rm sqrt(2)} = \beta\sqrt{c^2\boldsymbol\pi^2
- e c \hbar\,\boldsymbol\Sigma\!\cdot\!\mathbf{B}
+ (m c^2 + \beta V)^2}$
and yields
\begin{equation}
\label{eq:II3-clean}
i\hbar\,\frac{\partial\Psi}{\partial\tau} \;=\;
\left\{\frac{\boldsymbol\pi^2}{2 m}
 + \beta V
 + m c^2
 - \frac{e\hbar\,\boldsymbol\Sigma\!\cdot\!\mathbf{B}}{2 m c}
 + \frac{V^2}{2 m c^2}\right\}\Psi.
\end{equation}

\paragraph{Why this is structurally clean.}
Unlike~(II.1) and~(II.2), Eq.~(II.3) is scalar-only at the operator level
(apart from the trivial $\beta V$ term). Three commutations together
($[\beta, \boldsymbol\pi] = [\beta, \boldsymbol\Sigma] = [\beta, V] = 0$,
plus $\beta^2 = 1$) collapse the radicand without leaving operator-valued
cross terms. Eq.~(II.3) reads exactly like a Pauli-style wave equation with
a single relativistic correction $V^2/(2 m c^2)$, and in the
$V \to 0$ limit reduces cleanly to the Pauli equation plus a rest-energy
constant. \wolframcheck{(II.3) verified 2026-05-11; see
\S\ref{sec:eq-bs-s3}.}

\subsection{Dirac eigenvalue problem and the spinor split}

Section III recasts the Dirac equation as an eigenvalue problem with
$\Psi = [\psi_1, \psi_2]^t$:
\begin{align}
\label{eq:III1}
(\lambda - V - m c^2)\,\psi_1 &= c\,(\boldsymbol\sigma\!\cdot\!\boldsymbol\pi)\,\psi_2,\\
(\lambda - V + m c^2)\,\psi_2 &= c\,(\boldsymbol\sigma\!\cdot\!\boldsymbol\pi)\,\psi_1.
\nonumber
\end{align}
Solving the second for $\psi_2$:
\begin{equation}
\label{eq:III2}
\psi_2 \;=\; c\,\bigl[\,\lambda - V_0 + m c^2\,\bigr]^{-1}\,
            (\boldsymbol\sigma\!\cdot\!\boldsymbol\pi)\,\psi_1
\qquad\text{(III.2).}
\end{equation}

\subsection{From $H_D^2$ to the cutoff approximation}

Writing $H_D = H_0 + V_0$ with $H_0 = c\,\boldsymbol\alpha\!\cdot\!\boldsymbol\pi
+ \beta m c^2$ and using $\{\boldsymbol\alpha, \beta\} = 0$ together with the
Dirac identity
$(\boldsymbol\alpha\!\cdot\!\boldsymbol\pi)^2 = \boldsymbol\pi^2
- (e\hbar/c)\,\boldsymbol\Sigma\!\cdot\!\mathbf{B}$:
\begin{equation}
H_D^2 \;=\; H_0^2 + V_0^2 + \{H_0, V_0\},\qquad
H_0^2 \;=\; c^2 \boldsymbol\pi^2 - e \hbar c\,\boldsymbol\Sigma\!\cdot\!\mathbf{B}
+ m^2 c^4.
\end{equation}
Divide by $2 m c^2$, add $m c^2/2$, define
$V := \{H_0, V_0\}/(2 m c^2)$:
\begin{equation}
\label{eq:III3-KD}
K_D \;=\; \frac{\boldsymbol\pi^2}{2 m}
        + V
        - \frac{e\hbar\,\boldsymbol\Sigma\!\cdot\!\mathbf{B}}{2 m c}
        + m c^2
        + \frac{V_0^2}{2 m c^2}\qquad\text{(III.3).}
\end{equation}

\paragraph{Eq.~(III.4) --- $K_D$ with the shorthand $V$ expanded.}
Use $[\boldsymbol\alpha\!\cdot\!\boldsymbol\pi, V_0] = -i\hbar\,\boldsymbol\alpha\!\cdot\!\nabla V_0$
and $\{\beta, V_0\} = 2\beta V_0$ to expand
$V = \{H_0, V_0\}/(2 m c^2)$:
\begin{equation}
\label{eq:III4}
E\,\Psi \;=\; \Bigg\{
   \frac{\boldsymbol\pi^2}{2 m}
 + \beta V_0
 + m c^2
 - \frac{e\hbar\,\boldsymbol\Sigma\!\cdot\!\mathbf{B}}{2 m c}
 + \frac{V_0\,\boldsymbol\alpha\!\cdot\!\boldsymbol\pi}{m c}
 - \frac{i\hbar\,\boldsymbol\alpha\!\cdot\!\nabla V_0}{2 m c}
 + \frac{V_0^2}{2 m c^2}\,
\Bigg\}\Psi.
\end{equation}
\wolframcheck{Same expansion as (II.1) with $V \to V_0$; residual $0$
(DRQM~I~\S III.4, 2026-05-11).}

\paragraph{Eqs.~(III.5)--(III.6) --- upper/lower spinor split.}
Block-decompose~\eqref{eq:III4} into separate $\psi_1$ and $\psi_2$
equations. The $\psi_2$ equation is solved as~\eqref{eq:III2}. Substituting
into the $\psi_1$ equation produces a non-trivial chain-rule term because
$\boldsymbol\sigma\!\cdot\!\boldsymbol\pi$ does not commute with
$1/(\lambda - V_0 + m c^2)$:
\begin{equation}
\label{eq:III6-chain}
(\boldsymbol\sigma\!\cdot\!\boldsymbol\pi)\!\left[
\frac{(\boldsymbol\sigma\!\cdot\!\boldsymbol\pi)\,\psi_1}
     {\lambda - V_0 + m c^2}\right]
\;=\;
\frac{(\boldsymbol\sigma\!\cdot\!\boldsymbol\pi)^2\,\psi_1}
     {\lambda - V_0 + m c^2}
\;+\;
\frac{-i\hbar\,\boldsymbol\sigma\!\cdot\!\nabla V_0}
     {(\lambda - V_0 + m c^2)^2}\,
(\boldsymbol\sigma\!\cdot\!\boldsymbol\pi)\,\psi_1.
\end{equation}
The second piece of~\eqref{eq:III6-chain} (using
$\mathbf{p} = -i\hbar\nabla$, so $\boldsymbol\sigma\!\cdot\!\mathbf{p}\,V_0
= -i\hbar\,\boldsymbol\sigma\!\cdot\!\nabla V_0$) is the source of the
$(V_0/m)\,(\boldsymbol\sigma\!\cdot\!\mathbf{p}V_0)
(\boldsymbol\sigma\!\cdot\!\boldsymbol\pi)/(\lambda - V_0 + m c^2)^2\,\psi$
term in DRQM~I line~261. The result is the $\psi_1$-only eigenvalue
equation (III.6).

\paragraph{Eq.~(III.7) --- cutoff approximation.}
Approximating $\lambda - V_0 + m c^2 \approx 2 m c^2 + e^2/r$ for the
Coulomb potential $V_0 = -e^2/r$ (binding energy $\sim 10^{-5}\,m c^2$ is
dropped) and using $r_0 = e^2/(m c^2)$:
\begin{equation}
\label{eq:III7}
\psi_2 \;\approx\; \frac{c\,(\boldsymbol\sigma\!\cdot\!\boldsymbol\pi)}
                       {2 m c^2\!\left(1 + r_0/(2 r)\right)}\,\psi_1.
\end{equation}

\paragraph{Eq.~(III.8) --- approximated eigenvalue equation.}
Substituting~\eqref{eq:III7} into the $\psi_1$~equation from~(III.6) ---
i.e.~replacing every occurrence of $(\lambda - V_0 + m c^2)$ in (III.6) by
$2 m c^2\bigl(1 + r_0/(2 r)\bigr)$:
\begin{equation}
\label{eq:III8}
\begin{aligned}
\Bigl\{\frac{\boldsymbol\pi^2}{2m} + V_0 + \frac{V_0^2}{2 m c^2}
 - \frac{e\hbar\,\boldsymbol\sigma\!\cdot\!\mathbf{B}}{2 m c}
 + \frac{V_0\,(\boldsymbol\sigma\!\cdot\!\boldsymbol\pi)^2}{2 m^2 c^2\!\left(1 + r_0/(2r)\right)}
   \\[2pt]
{}- \frac{V_0\,(\boldsymbol\sigma\!\cdot\!\mathbf{p}\,V_0)\,(\boldsymbol\sigma\!\cdot\!\boldsymbol\pi)/m}
       {[\,2 m c^2(1 + r_0/(2r))\,]^2}
\Bigr\}\,\psi_1 \;=\; (E - m c^2)\,\psi_1.
\end{aligned}
\end{equation}
This is the approximated $\psi_1$-only eigenvalue equation that DRQM~I
(III.18)--(III.20) then collect into three identifiable Schr\"odinger-beyond
contributions.

\subsection{Eqs.~(III.9)--(III.10) --- Pauli identity}

The general Pauli identity (Sakurai Eq.~3.2.39):
\begin{equation}
\label{eq:III9-pauli}
(\boldsymbol\sigma\!\cdot\!\mathbf{X})(\boldsymbol\sigma\!\cdot\!\mathbf{Y})
\;=\; \mathbf{X}\!\cdot\!\mathbf{Y}
   \;+\; i\,\boldsymbol\sigma\!\cdot\!(\mathbf{X}\times\mathbf{Y})
\qquad\text{(III.9).}
\end{equation}
Specialising to $\mathbf{X} = \mathbf{Y} = \boldsymbol\pi$ with
$\boldsymbol\pi = \mathbf{p} - e\mathbf{A}/c$ and the gauge-field commutator
$[\pi_i, \pi_j] = i e\hbar\,\epsilon_{ijk} B_k/c$ (so
$\boldsymbol\pi\times\boldsymbol\pi = (i e\hbar/c)\mathbf{B}$):
\begin{equation}
\label{eq:III10-sigmapi-sq}
(\boldsymbol\sigma\!\cdot\!\boldsymbol\pi)^2
\;=\; \boldsymbol\pi^2 \;-\; \frac{e\hbar}{c}\,\boldsymbol\sigma\!\cdot\!\mathbf{B}
\qquad\text{(III.10).}
\end{equation}
\wolframcheck{\texttt{Simplify}$[\boldsymbol\pi^2 + i\cdot(i e\hbar/c)\boldsymbol\sigma\!\cdot\!\mathbf{B}]$
returns the (III.10) RHS (DRQM~I~\S III.10, verified).}

\subsection{Eqs.~(III.11)--(III.17) --- spherical-coordinate algebra}

Expand $(-i\hbar\,\boldsymbol\sigma\!\cdot\!\nabla V_0)(\boldsymbol\sigma\!\cdot\!\boldsymbol\pi)$
in spherical polar coordinates with the Coulomb scalar potential
$V_0 = -e^2/r$ and the proton dipole vector potential
$\mathbf{A} = (\boldsymbol\mu_p \times \mathbf{r})/r^3$. DRQM~I records the
chain across seven equations (III.11)--(III.17); the three load-bearing
results that the chain delivers (DRQM~I lines 329, 333, 338) are:
\begin{align}
\label{eq:III-spherical-1}
-i\hbar\,\nabla V_0\!\cdot\!\mathbf{p}
\;&=\; \frac{e^2 \hbar^2}{r^2}\,\frac{\partial}{\partial r},
&& \text{(scalar piece; $\sim$ velocity-radial coupling)} \\[2pt]
\label{eq:III-spherical-2}
\hbar\,\boldsymbol\sigma\!\cdot\!(\nabla V_0 \times \mathbf{p})
\;&=\; -\frac{e^2 \hbar}{r^3}\,\boldsymbol\sigma\!\cdot\!\mathbf{L},
&& \text{(spin--orbit emergence; }\mathbf{L} = \mathbf{r}\times\mathbf{p}\text{)} \\[2pt]
\label{eq:III-spherical-3}
-\frac{e\hbar}{c}\,\boldsymbol\sigma\!\cdot\!(\nabla V_0 \times \mathbf{A})
\;&=\; -\frac{e\hbar}{c}\cdot\frac{e^2}{r^4}\,
       2\,\mu_p\,|\mathbf{s}_p|\,\sin\theta\,
       (\boldsymbol\sigma\!\cdot\!\mathbf{e}_\theta).
&& \text{(electron-nuclear cross term)}
\end{align}
\eqref{eq:III-spherical-1} is the radial scalar piece (contributes to
relativistic kinetic + Darwin once the Pauli/FW reduction is taken);
\eqref{eq:III-spherical-2} is the standard spin--orbit emergence
$\boldsymbol\sigma\!\cdot\!\mathbf{L}/r^3$;
\eqref{eq:III-spherical-3} is the electron--proton dipole--dipole
contribution with explicit angular dependence
$\sin\theta\,(\boldsymbol\sigma\!\cdot\!\mathbf{e}_\theta)$.

\paragraph{Note on intermediate (III.11)--(III.16).}
DRQM~I records the spherical-coordinate chain across seven numbered
equations (III.11)--(III.17), but only the three load-bearing results
above carry forward to (III.18)--(III.20); the intermediate equations
(III.11)--(III.16) are the standard vector-calculus identities for
$\nabla(1/r)$, $\nabla V_0\!\cdot\!\mathbf{p}$, the cross-product
$\nabla V_0 \times \mathbf{p}$, and the same with $\mathbf{p}\to\mathbf{A}$
(proton vector potential), each evaluated in $(\hat r, \hat\theta, \hat\phi)$
basis. They appear in DRQM~I lines~297--339; the verification doc summary
records them as ``standard spherical-coordinate vector calculus, algebra
is mechanical.''

\subsection{The three ``new'' terms (III.18)--(III.20)}

Collecting~\eqref{eq:III-spherical-1}--\eqref{eq:III-spherical-3} together
with the Pauli-identity collapse~\eqref{eq:III10-sigmapi-sq} in the
approximated eigenvalue equation~\eqref{eq:III8}, three contributions
appear that are absent from the non-relativistic Schr\"odinger problem:
\begin{align}
\label{eq:III18-spin-mag}
\text{(III.18)}\quad
&-\Bigl[\,1 - \frac{4\,r_0}{2 r + r_0}\,\Bigr]
   \frac{e\hbar\,\boldsymbol\sigma\!\cdot\!\mathbf{B}}{2 m c}
   \qquad\textit{(spin--magnetic, cutoff-modified)} \\
\label{eq:III19-hyperfine-square}
\text{(III.19)}\quad
&\,2 r_0\,\mu_p^2\,|\mathbf{s}_p|^2\!
  \left[\,1 - \frac{4 e \hbar^2}{m c\,(2 r + r_0)}\,\right]
  \frac{\sin^2\theta}{r^4}
   \qquad\textit{(square of nuclear-dipole field)} \\
\label{eq:III20-hyperfine-cross}
\text{(III.20)}\quad
&-\frac{2 e\,r_0\,\hbar\,\mu_p\,|\mathbf{s}_p|}
       {m c\,(2 r + r_0)}\!
   \left[\,1 + \frac{4\,r_0}{2 r + r_0}\,\right]
   \frac{\sin\theta\,(\boldsymbol\sigma\!\cdot\!\mathbf{e}_\theta)}{r^3}
   \qquad\textit{(electron-nuclear cross term)}
\end{align}
\eqref{eq:III18-spin-mag} is the source of the modified $g$-factor;
\eqref{eq:III19-hyperfine-square}--\eqref{eq:III20-hyperfine-cross} are the
hyperfine pieces.

\subsection{The $g$-factor formula (III.21)--(III.22)}

Reading off the Zeeman coefficient from~\eqref{eq:III18-spin-mag} with
$\mathbf{s} = \hbar\boldsymbol\sigma/2$:
\begin{equation}
\label{eq:III21-anomalous-Ham}
H_a \;=\; 2\Bigl[\,1 - \frac{4\,r_0}{2 r + r_0}\,\Bigr]\,
         \mu_B\,\mathbf{s}\!\cdot\!\mathbf{B}\qquad\text{(III.21)},
\end{equation}
identifying the gyromagnetic ratio
\begin{equation}
\label{eq:III22-gr}
\boxed{\;g_r(x) \;=\; 2\Bigl[\,1 - \frac{4}{2 x + 1}\,\Bigr],
\qquad x \equiv r / r_0\;}\qquad\text{(III.22)}.
\end{equation}
Limit checks: $g_r(1/2) = -2$ (tree-Dirac), $g_r(x\to 0) = -6$.
For muon and proton, Eq.~(III.23) writes the analogue with
$r_0^\mu, r_0^p$ but leaves the cutoff values
$r_\mu, r_p$ as separate, numerically-unspecified parameters.

\paragraph{The closed-form algebraic inversion (\S III.D-extension).}
Inverting~\eqref{eq:III22-gr} for $g_r(x) = -2(1 + a_e)$ with $a_e$ the
electron anomalous magnetic moment, and solving algebraically:
\begin{equation}
\label{eq:rE-closed-form}
\boxed{\;\frac{r_e}{r_0} \;=\; \frac{2 - a_e}{2\,(2 + a_e)}\;}
\end{equation}
With CODATA-full $a_e^{\rm expt}$ this reproduces the triangulated
$r_e/r_0 = 0.499\,420\,509\,9128\dots$ to $3.45\times 10^{-13}$ (within the
precision floor $\sigma_r = 2.5\times 10^{-13}$); see
\S\ref{sec:eq-self-energy} and \S\ref{sec:z1-triangulation}. The
honest framing of~\eqref{eq:rE-closed-form} is
\emph{reproduction-by-inheritance under hypothesis~(ii)} (QED supplies
$a_e(\alpha)$; the framework supplies the cutoff--anomaly identification +
the algebraic inversion). An unconditional \textbf{Pass} requires
hypothesis-(i); see issue~\#75.

\subsection{From $g_r(r)$ to each observable}
\label{sec:from-gr-to-observables}

The campaign's six $g_s$-dependent precision observables each have a
prediction of the form
\begin{equation}
\label{eq:gs-n-textbook}
\mathcal{O}_i^{\rm framework}
\;=\; \left(\frac{g_s}{-2}\right)^{n_i}\!\!\!\cdot\;\mathcal{O}_i^{\rm textbook\,leading}
\qquad n_i \in \{1, 2\}
\end{equation}
with $g_s = g_r(r_e/r_0)$ from~\eqref{eq:III22-gr}, evaluated at the
triangulated cutoff. The $n_i$ exponents are determined by the operator
structure of each observable; the textbook leading term carries the rest
of the physics (kinematics, multipole structure, Bethe logs).

\subsubsection{Electron $g_s$ ($n=1$, directly from $g_r$)}

Set $r = r_e$ in~\eqref{eq:III22-gr}: the framework's prediction is
$g_s = g_r(r_e/r_0)$ directly. The triangulated $r_e/r_0$ is by construction
the value at which $g_s = -2.00231930\dots$ matches CODATA-2018, so the
match is exact by Pass-B fit construction
(\S\ref{sec:eq-triangulation}).

\subsubsection{Hydrogen $2P_{3/2}{-}2P_{1/2}$ fine structure ($n_{\rm FS}=1$)}

The Sommerfeld--Dirac formula gives the leading fine-structure splitting
\begin{equation}
\label{eq:fs-leading}
\Delta E_{\rm FS}^{\rm leading}(2P_{3/2}{-}2P_{1/2})
\;=\; \frac{m_e c^2\,(Z\alpha)^4}{32}\,(\text{angular factor})
\;\approx\; 10\,949\;\text{MHz at }Z=1.
\end{equation}
Spin--orbit is linear in $g_s$ (it enters via the $\mathbf{L}\!\cdot\!\mathbf{S}$
Lande factor in the Pauli/FW reduction of the dual Dirac equation,
BS~\S14.1), so $n_{\rm FS} = 1$:
\begin{equation}
\Delta E_{\rm FS}^{\rm framework}
\;=\; \frac{g_s}{-2}\,\Delta E_{\rm FS}^{\rm leading}
\;\approx\; 1.00115965 \cdot 10\,949
\;\approx\; 10\,962\;\text{MHz},
\end{equation}
with the framework anomalous correction $+12.7$~MHz on top of the leading
$10\,949$~MHz. Remaining $\sim 7$ MHz vs.~measured $10\,969.13(10)$ MHz is
recoil + two-loop, at the Bethe-estimate precision floor
(see \S\ref{sec:eq-li2plus-spectroscopy} for the $Z^4$ Li$^{2+}$ extension).

\subsubsection{Hydrogen 1S hyperfine (21-cm line, $n_{\rm HFS}=1$)}

The Fermi (1930) contact term gives the leading 1S hyperfine Hamiltonian
\begin{equation}
\label{eq:hfs-fermi}
H_{\rm HF}
\;=\; \frac{8\pi}{3}\,g_p\,\mu_N\,g_s\,\mu_B\,
       \delta^3(\mathbf{r})\,\mathbf{I}\!\cdot\!\mathbf{S},
\end{equation}
with $\langle 1S | \delta^3(\mathbf{r}) | 1S \rangle = 1/(\pi a_0^3)$,
producing the splitting
\begin{equation}
\Delta E_{\rm HF}^{\rm leading}(1S_{1/2})
\;=\; \frac{4}{3}\,g_p\,\frac{m_e}{M_p}\,\alpha^4\,m_e c^2
\;\approx\; 1\,418.4\;\text{MHz}\qquad(\text{at }g_s = -2).
\end{equation}
Hyperfine is \emph{linear} in $g_s$ ($n_{\rm HFS} = 1$, since the contact
term carries a single electron spin):
\begin{equation}
\Delta E_{\rm HF}^{\rm framework}
\;=\; \frac{g_s}{-2}\,\Delta E_{\rm HF}^{\rm leading}
\;\approx\; 1.00115965\cdot 1\,418.4
\;\approx\; 1\,420.04\;\text{MHz}.
\end{equation}
Residual vs.~measured $1\,420.405\,751\,768(2)$ MHz is $\sim 0.4$ MHz ---
sub-leading QED ($\alpha/\pi$ + recoil + nuclear structure),
$r_e$-independent, at the Bethe-estimate floor (BS-\S22.1).

\subsubsection{Lamb shift $2S_{1/2}{-}2P_{1/2}$ --- formulation-independent ($n = 0$)}

The Lamb shift escapes the $(g_s/{-}2)^n$ family. Bethe's (1947)
mass-renormalisation argument writes the self-energy shift
\begin{equation}
\label{eq:bethe-1947}
\Delta E_n^{\rm SE}
\;=\; \frac{2\alpha}{3\pi m^2 c^2}
   \sum_m |\langle m | \mathbf{p} | n \rangle|^2\,
         (E_m - E_n)\,
         \ln\!\frac{K}{|E_m - E_n|},
\end{equation}
where $K$ is a non-relativistic UV cutoff. Under the proper-time
formulation:
\begin{itemize}[topsep=2pt,itemsep=0pt]
  \item matrix elements $|\langle m | \mathbf{p} | n\rangle|^2$ are
        formulation-independent (BS-\S8);
  \item energy denominators $E_m - E_n$ are formulation-independent
        (BS-\S4);
  \item the Bethe-log structure $\ln(K/|E_m - E_n|)$ is non-relativistic in
        origin and inherits the same UV cutoff;
  \item the mass-renormalisation subtraction is formulation-independent at
        the order Bethe needed.
\end{itemize}
Therefore $\Delta E_{2S}^{\rm SE,\,proper-time} \approx 1\,040$ MHz is
\emph{identical} to the standard Bethe-estimate prediction at the precision
the route delivers. Adding the vacuum-polarisation shift ($-27$~MHz at 2S,
BS-\S21) and recombining gives $\sim 1\,013$ MHz at the Bethe-estimate
floor vs.~measured $1\,057.845(9)$ MHz; the $\sim 42$ MHz residual is the
gap to the full one-loop QED calculation, which the campaign does not
produce (BS-\S20).
$g_s$ does not enter the leading Bethe-log contribution
($\mathbf{p}\!\cdot\!\mathbf{p}$, not $\mathbf{s}\!\cdot\!\mathbf{B}$), so
$n_{\rm Lamb} = 0$ in~\eqref{eq:gs-n-textbook}.

\subsubsection{Helium ${}^3P_J$ fine structure ($n_{\rm He} = 1$)}

Two-electron spin--orbit + spin--spin recombination in the helium triplet
2P states produces the ${}^3P_0$--${}^3P_1$ interval. The leading
Bethe--Salpeter (1957) treatment gives a $(Z\alpha)^4$ scale modulated by
two-body recoil + correlation factors; spin--orbit is linear in $g_s$, so
$n_{\rm He} = 1$. Predicted framework value
$\sim 29\,616.95$ MHz vs.~measured $29\,616.952$ MHz at the kHz level,
provisional residual after spin--spin/spin--orbit recombination (BS-\S72).

\subsubsection{Positronium ortho--para ($n_{\rm Ps} = 2$)}

The $1{}^3S_1$--$1{}^1S_0$ ortho--para splitting in positronium ($e^+ e^-$)
involves a contact term \emph{quadratic} in the electron magnetic moment
(both leptons carry an anomalous moment that enters via the
two-fermion spin--spin coupling), giving $n_{\rm Ps} = 2$:
\begin{equation}
\Delta E_{\rm Ps}^{\rm framework}
\;=\; \left(\frac{g_s}{-2}\right)^{2}\!\cdot\Delta E_{\rm Ps}^{\rm textbook\,leading}
\;\approx\; (1.00115965)^2 \cdot 203\,389\;\text{MHz}.
\end{equation}
Per \texttt{10\_CrossComparison.md}~\S2, agreement is at Bethe precision;
positronium-specific QED (the annihilation contribution) is provisional in
the campaign (BS-\S80).

\subsubsection{Muonium hyperfine ($n_\mu = 1$)}

Muonium ($\mu^+ e^-$) replaces the proton in~\eqref{eq:hfs-fermi} with an
anti-muon. The Fermi-contact framework carries through with proton
$\to$ anti-muon kinematics; spin--orbit / Fermi contact is linear in
$g_s$, $n_\mu = 1$:
\begin{equation}
\Delta E_{\mu\rm{HFS}}^{\rm framework}
\;\approx\; \frac{g_s}{-2}\cdot 4\,458.2\;\text{MHz}
\;\approx\; 4\,463.4\;\text{MHz},
\end{equation}
vs.~measured $4\,463.302\,776(51)$ MHz: $\sim 0.1$ MHz residual at the
sub-leading-QED floor (BS-\S80).

\subsection{Summary of $n_i$ exponents and observable map}

\begin{center}
\renewcommand{\arraystretch}{1.05}
\begin{tabular}{lll}
\toprule
Observable                              & $n_i$ in $(g_s/{-}2)^{n_i}$ & Operator structure \\
\midrule
electron $g_s$                          & 1 (direct)   & $g_r(r_e/r_0)$ itself, Eq.~\eqref{eq:III22-gr} \\
H $2P_{3/2}{-}2P_{1/2}$                 & 1            & spin--orbit, $\mathbf{L}\!\cdot\!\mathbf{S}$ \\
H 1S hyperfine                          & 1            & Fermi contact, $\delta^3(\mathbf{r})\,\mathbf{I}\!\cdot\!\mathbf{S}$ \\
He ${}^3P_0{-}^3P_1$                    & 1            & two-body spin--orbit (post-recombination) \\
positronium ortho--para                 & 2            & two-fermion spin--spin contact (both leptons) \\
muonium hyperfine                       & 1            & Fermi contact, $e^- \mu^+$ \\
H Lamb shift $2S{-}2P$                  & 0            & formulation-independent self-energy (Bethe 1947) \\
\bottomrule
\end{tabular}
\end{center}

\paragraph{What the chain establishes (and does not).}
Once~\eqref{eq:III22-gr} is in hand and the triangulated $r_e/r_0$ is
fixed by Pass-B joint fit (\S\ref{sec:z1-triangulation}), each
$g_s$-dependent observable evaluates as $(g_s/{-2})^{n_i}$ $\times$ its textbook
leading term --- a single back-fit applied six times rather than six
independent corroborations
(\texttt{10\_CrossComparison.md}~\S2; \S\ref{sec:findings-per-observable}).
The Lamb shift escapes the family entirely because the leading Bethe-log
self-energy is $g_s$-symmetric. The chain's content is the (III.22)
cutoff--anomaly identification; the empirical $a_e^{\rm bound}(Z\alpha)$
curve across the Z-scan is QED's, per the
\hyperref[sec:branch-a-vs-branch-c]{Branch~A vs.~Branch~C} reading.

\authorreview{This section transcribes the derivation chain you already
laid out in DRQM~I \S\S II--III, with every typeset numbered equation from
(II.1) through (III.22) surfaced for fine-grain review per your
2026-05-30 instruction. The intermediate (III.11)--(III.16)
spherical-coordinate-algebra equations are noted as ``mechanical
vector-calculus'' per the verification doc's own characterisation and not
re-derived here; their three load-bearing
outputs~\eqref{eq:III-spherical-1}--\eqref{eq:III-spherical-3} are
displayed. The per-observable $n_i$ summary table is a synthesis move
(drawn from the BS section docs but consolidated here for the first time in
one place) and is endorsed per your 2026-05-30 sign-off.}
